% Part: computability
% Chapter: computability-theory
% Section: k-1

\documentclass[../../../include/open-logic-section]{subfiles}

\begin{document}

\olfileid{cmp}{thy}{k1}
\olsection{হ্রাসযোগ্যতার একটি উদাহরণ}

\olref[ppr]{prop:reduce}-এর একটি প্রয়োগ বিবেচনা করা যাক।

\begin{prop}
\ollabel{prop:k1}
ধরা যাক,
\[
K_1 = \Setabs{e}{\cfind{e}(0) \fdefined}.
\]
তাহলে $K_1$ গণনসাধ্যভাবে তালিকায়নযোগ্য, কিন্তু গণনসাধ্য নয়।
\end{prop}

\begin{proof}
যেহেতু $K_1 = \Setabs{e}{\lexists[s][T(e,0,s)]}$, তাই
\olref[eqc]{thm:exists-char} অনুসারে $K_1$ গণনসাধ্যভাবে তালিকায়নযোগ্য।

$K_1$ গণনসাধ্য নয় দেখাতে আমরা দেখাব যে $K_0$ তার মধ্যে হ্রাসযোগ্য।

\begin{explain}
এটি একটু জটিল, কারণ $K_1$ ব্যবহার করে আমরা কেবল একটি নির্দিষ্ট নিবেশ~$0$
দিয়ে শুরু হওয়া গণনা সম্পর্কে প্রশ্ন করতে পারি। ধরা যাক, তোমার এমন একজন
বুদ্ধিমান বন্ধু আছে যে এই ধরনের প্রশ্নের উত্তর দিতে পারে (এমন বন্ধুদের
“ওরাকল” বলা হয়)। এরপর কেউ এসে জিজ্ঞাসা করল, $\tuple{e,x}$, $K_0$-তে
আছে কি না, অর্থাৎ যন্ত্র~$e$ নিবেশ~$x$-এ থামে কি না। তুমি আরেকটি
যন্ত্র~$e_x$ বানাতে পারো, যা \emph{যেকোনো} নিবেশে সেই নিবেশটি উপেক্ষা
করে তার বদলে নিবেশ~$x$-এ যন্ত্র~$e$ চালায়। তাহলে যন্ত্র~$e$ নিবেশ~$x$-এ
থামে কি না—এই প্রশ্নটি যন্ত্র~$e_x$ নিবেশ~$0$-তে (বা অন্য যেকোনো নিবেশে)
থামে কি না—এই প্রশ্নটির সমতুল্য। তাই নতুন যন্ত্র~$e_x$ নিবেশ~$0$-তে
থামে কি না তুমি বন্ধুকে জিজ্ঞাসা করো; বদলানো প্রশ্নটির উত্তরে মূল
প্রশ্নটির উত্তর পাওয়া যায়। এটি $K_0$ থেকে~$K_1$-তে কাঙ্ক্ষিত
হ্রাসকরণটি দেয়।
\end{explain}

সার্বজনীন আংশিক গণনসাধ্য অপেক্ষক ব্যবহার করে $f$-কে নিচে সংজ্ঞায়িত
3-যুক্তিক অপেক্ষক ধরি:
\[
f(x,y,z) \simeq \cfind{x}(y).
\]
লক্ষ করো, $f$ তার তৃতীয় নিবেশটি সম্পূর্ণ উপেক্ষা করে। এমন একটি সূচক~$e$
বেছে নিই যাতে $f = \cfind{e}[3]$; ফলে
\[
\cfind{e}[3](x,y,z) \simeq \cfind{x}(y).
\]
$s$-$m$-$n$ উপপাদ্য অনুসারে এমন একটি অপেক্ষক $s(e,x,y)$ আছে, যাতে
প্রত্যেক~$z$-এর জন্য
\begin{align*}
\cfind{s(e,x,y)}(z) & \simeq \cfind{e}[3](x,y,z) \\
& \simeq \cfind{x}(y).
\end{align*}

\begin{explain}
উপরের অনানুষ্ঠানিক যুক্তির ভাষায়, $s(e,x,y)$ হলো সেই যন্ত্রটির একটি
সূচক, যা যেকোনো নিবেশ~$z$-এর জন্য সেই নিবেশটি উপেক্ষা করে
$\cfind{x}(y)$ গণনা করে।
\end{explain}

বিশেষ করে,
\[
\cfind{s(e,x,y)}(0) \fdefined \quad \text{যদি এবং কেবল যদি} \quad
\cfind{x}(y) \fdefined.
\]
অন্যভাবে বললে, $\tuple{x, y} \in K_0$ যদি এবং কেবল যদি
$s(e,x,y) \in K_1$। তাই নিচে সংজ্ঞায়িত অপেক্ষক~$g$
\[
g(w) = s(e,(w)_0,(w)_1)
\]
হলো $K_0$ থেকে~$K_1$-তে একটি হ্রাসকরণ।
\end{proof}

\end{document}
